\clearpage
%\fancyfood[R]{\thepage}
%\renewcommand{\headrulewidth}{0.5pt}
\chapter{Project general context}
\chaptermark{Project general context}
\pagestyle{fancy}
\fancyhf{}
\fancyhead[HL]{Chapter 1: Project general context}
\cfoot{\thepage}

\section{Introduction}
%__TODO: NEEDS UPDATE
\hspace{1cm}This Chapter is structured as follow: first, we'll start by giving a brief introduction to the academical institution and the host company where the graduation project was developed. After that, we'll be interested in presenting the didactic model and outlining its principal features and components. Finally, we'll wrap up this chapter by highlighting the system main purposes.
\section{ISIMM presentation}
\hspace{1cm}The Higher Institute of Computer Science and Mathematics of the University of Monastir (ISIMM) is created by the decree n° 1623 of July 09, 2002, is a scientific, public higher education establishment, placed under the supervision of the Ministry of Higher Education and Scientific Research.

ISIMM offers educational programs in the field of Science and Technology that focus on the following disciplines: Computer Science, Mathematics and Electronics\cite{0}.

%__ Hosting company presentation 
\section{Company overview}

\subsection{Presentation}
\hspace{1cm}
BWS (Be Wireless Solutions) is a Tunisian technology company specializing in wireless connectivity and IoT solutions. By combining advanced wireless technologies with reliable embedded systems, BWS provides innovative services that enable seamless communication between devices and networks. The company focuses on designing and integrating robust, scalable solutions for industrial and smart‑environment applications\cite{1}.

\subsection{Company industries}
\subsubsection{Wireless connectivity solutions}
\hspace{1cm}
A core activity of BWS is developing and deploying wireless connectivity solutions based on technologies such as Wi‑Fi, Bluetooth, LoRaWAN and other wireless protocols. These solutions are tailored to specific customer needs, including long‑range communication, low‑power operation, and secure data transmission. BWS ensures that its wireless systems are optimized for performance, reliability, and energy efficiency.
\subsubsection{IoT and embedded systems}
\hspace{1cm}
BWS designs and implements end‑to‑end IoT solutions using embedded systems and microcontrollers. The company handles hardware selection, firmware development, and communication protocols to create smart devices that can collect, process, and transmit data in real time. These solutions are widely used in industrial monitoring, environmental sensing, and smart building applications.
\subsubsection{Industrial and smart‑environment applications}
\hspace{1cm}
BWS applies its wireless and IoT expertise to industrial and smart‑environment use cases, such as remote monitoring, asset tracking, and automation. The company’s solutions help customers reduce operational costs, improve safety, and increase efficiency by providing accurate, real‑time data and remote control capabilities. BWS emphasizes robustness, scalability, and adaptability to different environments, from factories to smart cities.

%__ Project specification
\section{Project specification}
\subsection{Objectives}

\hspace{1cm}
The main goals of the project are:
\begin{itemize}
    \item Designing and developing the \textbf{IO BOX–BWS} platform for industrial signal acquisition, processing, and regulation.
    \item Acquiring \textbf{analog and digital signals} from sensors and industrial equipment in a reliable and real‑time manner.
    \item Processing and filtering the acquired data to improve signal quality and extract useful information.
    \item Implementing control and regulation functions to drive actuators and keep processes stable.
    \item Supporting several industrial communication protocols to connect with sensors, actuators, and external systems.
    \item Building a simple, modular, and maintainable system that can be reused in different industrial applications.
\end{itemize}

\subsection{Technical specifications}
\hspace{1cm}
The IO BOX–BWS platform must meet the following technical specifications:
\begin{itemize}
    \item \textbf{Signal acquisition:} The platform must acquire analog and digital signals from industrial sensors and equipment.
    \item \textbf{Signal processing:} The platform must apply basic filters and calculations (such as averaging and scaling) to the acquired signals.
    \item \textbf{Control and regulation:} The platform must support simple control laws (such as PID) to regulate process variables.
    \item \textbf{Communication:} The platform must communicate with sensors and actuators using standard industrial protocols.
    \item \textbf{External interface:} The platform must send processed data to external systems or supervision software.
    \item \textbf{Modularity:} The software must be organized in clear, independent modules to make it easy to modify and reuse.
\end{itemize}